% Options for packages loaded elsewhere
\PassOptionsToPackage{unicode}{hyperref}
\PassOptionsToPackage{hyphens}{url}
%
\documentclass[
  10pt,
  english,
]{article}
\usepackage{amsmath,amssymb}
\usepackage{lmodern}
\usepackage{iftex}
\ifPDFTeX
  \usepackage[T1]{fontenc}
  \usepackage[utf8]{inputenc}
  \usepackage{textcomp} % provide euro and other symbols
\else % if luatex or xetex
  \usepackage{unicode-math}
  \defaultfontfeatures{Scale=MatchLowercase}
  \defaultfontfeatures[\rmfamily]{Ligatures=TeX,Scale=1}
\fi
% Use upquote if available, for straight quotes in verbatim environments
\IfFileExists{upquote.sty}{\usepackage{upquote}}{}
\IfFileExists{microtype.sty}{% use microtype if available
  \usepackage[]{microtype}
  \UseMicrotypeSet[protrusion]{basicmath} % disable protrusion for tt fonts
}{}
\makeatletter
\@ifundefined{KOMAClassName}{% if non-KOMA class
  \IfFileExists{parskip.sty}{%
    \usepackage{parskip}
  }{% else
    \setlength{\parindent}{0pt}
    \setlength{\parskip}{6pt plus 2pt minus 1pt}}
}{% if KOMA class
  \KOMAoptions{parskip=half}}
\makeatother
\usepackage{xcolor}
\IfFileExists{xurl.sty}{\usepackage{xurl}}{} % add URL line breaks if available
\IfFileExists{bookmark.sty}{\usepackage{bookmark}}{\usepackage{hyperref}}
\hypersetup{
  pdftitle={Coastal Representation Learning by Combining Optical and Depth Foundation Models},
  pdflang={en},
  hidelinks,
  pdfcreator={LaTeX via pandoc}}
\urlstyle{same} % disable monospaced font for URLs
\usepackage[margin=0.85in]{geometry}
\usepackage{longtable,booktabs,array}
\usepackage{calc} % for calculating minipage widths
% Correct order of tables after \paragraph or \subparagraph
\usepackage{etoolbox}
\makeatletter
\patchcmd\longtable{\par}{\if@noskipsec\mbox{}\fi\par}{}{}
\makeatother
% Allow footnotes in longtable head/foot
\IfFileExists{footnotehyper.sty}{\usepackage{footnotehyper}}{\usepackage{footnote}}
\makesavenoteenv{longtable}
\usepackage{graphicx}
\makeatletter
\def\maxwidth{\ifdim\Gin@nat@width>\linewidth\linewidth\else\Gin@nat@width\fi}
\def\maxheight{\ifdim\Gin@nat@height>\textheight\textheight\else\Gin@nat@height\fi}
\makeatother
% Scale images if necessary, so that they will not overflow the page
% margins by default, and it is still possible to overwrite the defaults
% using explicit options in \includegraphics[width, height, ...]{}
\setkeys{Gin}{width=\maxwidth,height=\maxheight,keepaspectratio}
% Set default figure placement to htbp
\makeatletter
\def\fps@figure{htbp}
\makeatother
\setlength{\emergencystretch}{3em} % prevent overfull lines
\providecommand{\tightlist}{%
  \setlength{\itemsep}{0pt}\setlength{\parskip}{0pt}}
\setcounter{secnumdepth}{-\maxdimen} % remove section numbering
\ifXeTeX
  % Load polyglossia as late as possible: uses bidi with RTL langages (e.g. Hebrew, Arabic)
  \usepackage{polyglossia}
  \setmainlanguage[]{english}
\else
  \usepackage[main=english]{babel}
% get rid of language-specific shorthands (see #6817):
\let\LanguageShortHands\languageshorthands
\def\languageshorthands#1{}
\fi
\ifLuaTeX
  \usepackage{selnolig}  % disable illegal ligatures
\fi

\title{Coastal Representation Learning by Combining Optical and Depth
Foundation Models}
\author{}
\date{Working draft --- 3 November 2026}

\begin{document}
\maketitle

\hypertarget{abstract}{%
\section{Abstract}\label{abstract}}

Coastal prediction combines spectral appearance, spatial context and
geometric structure, but these signals are rarely provided by one
pretrained encoder. We study a simple dual-foundation architecture that
combines an optical Earth-observation model, OLMoEarth Large, with the
visual encoder of PromptDA, a depth foundation model. Both encoders
remain frozen: their spatial tokens are aligned and concatenated before
a task-specific linear probe. The evaluation uses four downstream
tasks---bathymetry regression and three segmentation benchmarks---under
a common validation-selected probing protocol. On 1,150 coastal regions
used for coastal adaptation experiments, the dual-foundation
representation improves MARIDA mIoU from 0.6129 to 0.6384 and MADOS mIoU
from 0.6480 to 0.6833 relative to OLMoEarth Large alone. Bathymetry
changes from 0.5771 to 0.5749 m RMSE, whereas LICS decreases from 0.9548
to 0.9431. We report OLMoEarth Nano adaptation and a task-adapted
SpectralMoE implementation as controls that separate the dual-foundation
result from additional coastal or target-task supervision. The results
support optical--depth feature complementarity on selected coastal
tasks, while the single-run design, larger concatenated probe and
unmatched interfaces limit claims of general superiority.

\hypertarget{introduction}{%
\section{1. Introduction}\label{introduction}}

Coastal prediction spans tasks with different information requirements.
Shoreline delineation depends on boundaries, marine surface
interpretation depends on spectral and contextual cues, and bathymetry
requires continuous geometric estimates. Optical Earth-observation
models are well matched to spectral variation, but a depth foundation
model can provide a complementary bias toward scene layout and relative
structure. The central question is whether these two pretrained
representations can be combined without retraining either encoder.

OLMoEarth provides a multimodal optical representation and a
frozen-feature evaluation recipe {[}1{]}. PromptDA extends the Depth
Anything line toward high-resolution metric-depth estimation {[}3{]};
here we use its visual encoder features, not its depth predictions or
depth prompts. Combining the two models is attractive because it adds a
geometric representation without requiring paired depth labels in the
coastal collection. The combination also creates practical questions
about spatial alignment, feature dimensionality and whether any gain
survives a frozen, task-independent evaluation.

We investigate the dual-foundation design through a local evaluation
study rather than new pretraining. We first establish ten
pretrained-model references under a common probing budget. We then
compare OLMoEarth Large alone with frozen OLMoEarth--PromptDA
concatenation on four tasks. OLMoEarth Nano coastal adaptation and a
separate transfer of the released SpectralMoE architecture {[}2{]}
provide controls for the effects of coastal supervision and target-task
adaptation. These controls are reported separately because they change
the representation before linear probing.

Our contribution is an explicit and reproducible baseline for combining
optical and depth foundation-model features in coastal remote sensing.
The method has three deliberately small operations: independent
preprocessing, spatial-token alignment and channel concatenation,
followed by a task-specific probe. The experiments show selected gains
on marine-surface segmentation while exposing a trade-off on shoreline
segmentation. We do not claim a new depth model, a new optical encoder
or a new fusion primitive; the scientific claim is narrower: a frozen
depth foundation model can complement an optical foundation model on
some coastal transfers under the tested protocol.

\hypertarget{data-and-evaluation-settings}{%
\section{2. Data and evaluation
settings}\label{data-and-evaluation-settings}}

\hypertarget{coastal-adaptation-collection}{%
\subsection{2.1. Coastal adaptation
collection}\label{coastal-adaptation-collection}}

We use 1,150 coastal regions selected from an existing collection for
complete monthly records in Sentinel-2, Landsat, and MODIS during 2024.
The split contains 911 training, 120 validation, and 119 held-out
regions. Months from a region retain the same split membership. Complete
monthly records do not imply that every pixel is valid; observation
masks remain necessary.

Sentinel-2 supplies 12 bands per time step. Landsat uses six scientific
bands, organized through the existing input adapter. MODIS supplies
seven bands on a native 5 by 5 grid, represented as 25 coarse tokens per
time step. The optical training footprint spans approximately 2,560
projected metres; the MODIS grid spans approximately 2,500 projected
metres. MODIS is not upsampled to simulate fine spatial detail. The
current experiments do not use Sentinel-1, although an archive audit
found all 12 monthly Sentinel-1 records for 1,027 of the regions.

Auxiliary supervision comprises shoreline labels from an existing JRC
2020--2024 aggregation, GLC-FCS30D land-cover labels from 2022, a
Copernicus GLO30 static surface-elevation product, and
tide-model-derived monthly targets. The elevation product is a digital
surface model, not underwater bathymetry, and its edition year does not
imply acquisition in 2024. Spatial targets are aligned to a common
origin and a 100 by 100 grid with 25.6 m spacing. Differences in source
dates mean that these targets are not fully synchronous observations.

The regions span approximately 47.08 degrees south to 51.93 degrees
north. Europe and western Asia contribute 275 regions, North America
245, Africa and the Middle East 232, South America 149, Oceania 146,
East Asia and the northern Pacific 57, and southern and southeastern
Asia 46. Thus, coverage is geographically broad but uneven. These
regions have not yet undergone a spatial-overlap audit against
downstream datasets or the foundation models' pretraining corpora.

\begin{longtable}[]{@{}lrr@{}}
\toprule
Water fraction within valid JRC pixels & All regions & Training
regions \\
\midrule
\endhead
0\% & 150 & 121 \\
More than 0\%, at most 25\% & 309 & 246 \\
More than 25\%, less than 75\% & 320 & 250 \\
At least 75\%, less than 100\% & 232 & 178 \\
100\% & 139 & 116 \\
Total & 1,150 & 911 \\
\bottomrule
\end{longtable}

Table 1. Composition of aligned coastal labels. Fractions describe valid
observed pixels, not the entire unmasked tile. All-water regions are not
necessarily exclusively marine. Original sampling-footprint fractions
differ from these cropped-label statistics.

\hypertarget{downstream-tasks}{%
\subsection{2.2. Downstream tasks}\label{downstream-tasks}}

We evaluate MagicBathyNet bathymetry regression, LICS shoreline
segmentation, MARIDA marine surface segmentation, and MADOS marine
surface segmentation. Bathymetry is measured by root mean squared error
(RMSE) in metres; segmentation is measured by mean intersection over
union (mIoU). Dataset names identify the local benchmark inputs; our
adaptations and splits should not be interpreted as a universal
reproduction of published dataset leaderboards.

\begin{longtable}[]{@{}lrrrl@{}}
\toprule
Benchmark & Train & Validation & Test & Metric \\
\midrule
\endhead
MagicBathyNet & 2,030 & 226 & 566 & RMSE, lower is better \\
LICS & 80 & 20 & 100 & mIoU, higher is better \\
MARIDA & 694 & 328 & 359 & mIoU, higher is better \\
MADOS & 2,657 & 1,186 & 1,311 & mIoU, higher is better \\
\bottomrule
\end{longtable}

Table 2. Sample counts in the local evaluation manifests. LICS uses the
first 80 samples of the local fine-tuning partition for training and the
remaining 20 for validation. MagicBathyNet uses the existing stored
split IDs; MARIDA and MADOS use the partitions exposed by their local
loaders. Canonical dataset citations, acquisition details, and release
identifiers must accompany a public release of these manifests.

\hypertarget{linear-probing}{%
\subsection{2.3. Linear probing}\label{linear-probing}}

We freeze the feature extractor and train a fresh linear probe for each
downstream task. Following the OLMoEarth evaluation design {[}1{]}, each
trial runs for 50 epochs. The learning-rate candidates are 0.0001,
0.0005, 0.001, 0.005, 0.01, 0.05, 0.1, and 0.5. AdamW uses weight decay
0.01, batch size 32, and five epochs of linear warmup followed by cosine
decay towards 0.00001. Probe initialization uses seed 42. Validation
selects the configuration after each 50-epoch trial; the selected
configuration is then evaluated on the test split. Test scores do not
select the learning rate or normalization.

For OLMoEarth, the input patch size is four. We compare the official
pretrained normalization statistics with benchmark statistics using the
same mean-and-standard-deviation scaling. Most benchmark statistics are
fitted on the training inputs; MADOS uses published statistics embedded
in the existing evaluation implementation. The current tasks use
single-time-step inputs. External models retain their native adapters
and input-size requirements, so identical optimization budgets do not
imply identical sensor information or preprocessing.

For segmentation, a linear layer maps each spatial feature to class
logits for its corresponding output patch; logits are resized to the
original target grid when required. Cross-entropy ignores label 255 and
does not use additional class weights. For regression, a single-channel
1 by 1 convolution predicts bathymetry and minimizes mean squared error
over valid pixels. Neither task uses the coastal multitask loss during
probing. Features are extracted with BF16 autocast and stored in float16
caches in the current implementation.

\hypertarget{opticaldepth-dual-foundation-method-and-controls}{%
\section{3. Optical--depth dual-foundation method and
controls}\label{opticaldepth-dual-foundation-method-and-controls}}

\hypertarget{official-checkpoint-references}{%
\subsection{3.1. Official checkpoint
references}\label{official-checkpoint-references}}

We evaluate official OLMoEarth Nano and Large checkpoints alongside
TerraMind, Prithvi EO-2.0, DOFA, SatMAE, CROMA, Scale-MAE, AnySat, and
TESSERA. OLMoEarth Nano and Large emit 128- and 1,024-dimensional
features, respectively. Several interfaces require sensor or band
adaptation. Scale-MAE uses RGB, whereas the other adapters can use
different spectral subsets. The present TESSERA experiment supplies a
single optical observation with the existing temporal-bin adaptation and
does not evaluate its full annual multimodal setting. Lower scores under
these interfaces should not be taken as upper bounds on the underlying
models' capabilities.

\hypertarget{olmoearthpromptda-feature-fusion}{%
\subsection{3.2. OLMoEarth--PromptDA feature
fusion}\label{olmoearthpromptda-feature-fusion}}

The primary method combines the official OLMoEarth Large optical encoder
with the official PromptDA ViT-L depth foundation model {[}3{]}.
OLMoEarth receives the optical bands through its native sensor adapter.
PromptDA receives an RGB rendering of the same observation through its
native image preprocessing. The two encoders are evaluated independently
and remain frozen throughout feature extraction and probing.

The OLMoEarth input is 256 by 256 with patch size four in the probing
protocol. PromptDA receives 448 by 448 RGB inputs with its native patch
size fourteen. Its 32 by 32 final normalized patch-token grid is
bilinearly aligned with the OLMoEarth spatial grid. We concatenate the
aligned channels:

\[
F_{\mathrm{dual}}=[F_{\mathrm{OLMoEarth}};\operatorname{Resize}(F_{\mathrm{PromptDA}})]\in\mathbb{R}^{H\times W\times 2048}.
\]

PromptDA contributes a representation learned for depth estimation; the
fusion does not use predicted depth, ground-truth depth, a depth prompt,
or a learned projection. The same two encoder weights are used for all
four tasks. A fresh linear probe is trained on the concatenated
features, so the reported dual-foundation result measures transfer of a
task-independent representation. The concatenated probe has twice the
feature width of the OLMoEarth-only probe, an acknowledged source of
additional capacity.

\hypertarget{coastal-adaptation-of-nano-as-a-control}{%
\subsection{3.3. Coastal adaptation of Nano as a
control}\label{coastal-adaptation-of-nano-as-a-control}}

The no-MoE baseline retains the official Nano transformer and adds a
native MODIS input adapter and coastal task heads. The training
objective is

\[
L_{\mathrm{coast}}=L_{\mathrm{shore}}+1.6L_{\mathrm{landcover}}+0.1L_{\mathrm{elevation}}+0.5L_{\mathrm{tide}}.
\]

Shoreline uses binary cross-entropy, land cover uses cross-entropy, and
elevation and tide use SmoothL1. Elevation targets use a signed
logarithmic transformation in the coastal training implementation.
Training runs for 40 epochs with batch size four and seed 42. The
official parameters are frozen for the first three epochs, then updated
at learning rate 0.000045; added modules use 0.0045. AdamW weight decay
is 0.05. A separate coastal validation objective selects the checkpoint;
it uses a 0.1 coefficient for mean tide and is not identical to the
training objective.

The Nano MoE variants use six experts with corrected Top-1 or Top-2
routing and physical token coordinates. Coordinates encode token centres
using each modality's pixel spacing and origin before common
normalization. This changes relative spatial representation; it does not
supply a new global location signal. Corrected Top-1 uses a
straight-through gate because softmax over a single selected logit is
constant. Load-balancing and expert-diversity penalties have
coefficients 0.01 and 0.001, with a three-epoch diversity warmup.

The unfrozen Top-1 and Top-2 runs use the coastal validation-selected
epochs 39 and 30. In a separate experiment, the official Nano parameters
remain frozen throughout training while added modules and heads are
updated. These runs use seed 43 and BF16; we report their fixed epoch-40
snapshots. Consequently, frozen versus unfrozen comparisons also differ
in seed, precision, and checkpoint selection and are not isolated causal
tests of freezing. Coastal training uses patch size eight, whereas
downstream probing uses patch size four. Downstream inputs do not
include paired MODIS, and the coastal task heads are discarded.

\hypertarget{target-task-adapted-spectralmoe-as-a-secondary-control}{%
\subsection{3.4. Target-task-adapted SpectralMoE as a secondary
control}\label{target-task-adapted-spectralmoe-as-a-secondary-control}}

We also adapt the released SpectralMoE architecture of Chen et
al.~{[}2{]}, replacing its visual foundation model with OLMoEarth Large
while retaining the PromptDA encoder and the released six-expert, Top-1
adapter design. This is an implementation adaptation of prior work, not
a newly introduced dual-gated MoE architecture. Visual inputs use 256 by
256 with patch eight, and PromptDA uses 448 by 448 with patch fourteen.
The interface aligns the spatial grids and averages the visual spectral
groups as implemented.

Each downstream task separately trains MoE and task-head parameters for
20 epochs with the two foundation encoders frozen. Segmentation uses the
adapted Mask2Former head; bathymetry uses a regression head and MSE. We
then discard the task head, freeze the adapted encoder and MoE, and
train a new linear probe. The representation has therefore already
received target-task supervision before probing. The adapter weights
differ by task, and these scores cannot be interpreted as those of a
single task-independent foundation model. Interface adaptations and the
retained author implementation also prevent claiming an exact
reproduction of SpectralMoE's original benchmark recipe.

\hypertarget{results}{%
\section{4. Results}\label{results}}

\begin{figure}
\centering
\includegraphics{transfer_comparison.png}
\caption{Figure 1. Task profiles of official Nano, coastally adapted
Nano without MoE, official OLMoEarth Large, and the frozen
OLMoEarth--PromptDA dual-foundation representation. Values are final
test metrics under the local linear-probing protocol; lower is better
for RMSE and higher for mIoU. Bars start at zero. Each bar represents
one run, without uncertainty estimates.}
\end{figure}

\hypertarget{official-models-provide-different-reference-points}{%
\subsection{4.1. Official models provide different reference
points}\label{official-models-provide-different-reference-points}}

No official checkpoint leads all four local tasks (Table 3). TerraMind
achieves the lowest bathymetry RMSE and highest MADOS mIoU in this
group. Prithvi leads LICS, whereas Nano leads MARIDA. The contrast
between Nano and Large on MARIDA cautions against treating encoder size
as a reliable predictor of performance under the current sensor
interfaces and probing configuration.

\begin{longtable}[]{@{}lrrrr@{}}
\toprule
Official checkpoint & Magic RMSE & LICS mIoU & MARIDA mIoU & MADOS
mIoU \\
\midrule
\endhead
OLMoEarth Nano & 0.633355 & 0.955118 & \textbf{0.647160} & 0.565515 \\
OLMoEarth Large & 0.577098 & 0.954755 & 0.612900 & 0.647969 \\
TerraMind Base & \textbf{0.529318} & 0.959422 & 0.499219 &
\textbf{0.716325} \\
Prithvi EO-2.0 300M & 0.605855 & \textbf{0.968702} & 0.493729 &
0.641364 \\
DOFA ViT-Base & 0.677503 & 0.911565 & 0.413663 & 0.549592 \\
SatMAE ViT-Base multispectral & 0.621436 & 0.955805 & 0.374972 &
0.617544 \\
CROMA Base & 1.397860 & 0.922307 & 0.198291 & 0.212326 \\
Scale-MAE ViT-Large & 0.664649 & 0.709480 & 0.174337 & 0.333993 \\
AnySat Base & 0.556970 & 0.442266 & 0.239017 & 0.222442 \\
TESSERA v2 Medium & 0.698643 & 0.932338 & 0.365159 & 0.500293 \\
\bottomrule
\end{longtable}

Table 3. Official checkpoints under local frozen probing. Bold marks the
best value in each column. Input adaptations differ; these are local
evaluation results rather than official reported model scores.

\hypertarget{coastal-adaptation-changes-the-task-profile}{%
\subsection{4.2. Coastal adaptation changes the task
profile}\label{coastal-adaptation-changes-the-task-profile}}

Coastal supervision improves the no-MoE Nano baseline on LICS by 1.46
percentage points but reduces its MARIDA and MADOS scores by 13.96 and
4.66 points, respectively (Table 4). Its bathymetry RMSE also increases.
The combination of a shoreline gain and marine-task losses is consistent
with specialization, but does not uniquely establish catastrophic
forgetting: altered input handling, checkpoint selection, and patch-size
transfer are alternative contributors.

Keeping the official encoder frozen changes this trade-off. Fixed
epoch-40 Top-1 adaptation yields bathymetry RMSE of 0.568042 m, a
10.31\% reduction relative to the official Nano reference, while
remaining below that reference on all three segmentation tasks. Top-2
does not consistently dominate Top-1. We therefore report these variants
as different task profiles, rather than identifying one as a generally
better coastal foundation model.

\begin{longtable}[]{@{}
  >{\raggedright\arraybackslash}p{(\columnwidth - 8\tabcolsep) * \real{0.16}}
  >{\raggedleft\arraybackslash}p{(\columnwidth - 8\tabcolsep) * \real{0.21}}
  >{\raggedleft\arraybackslash}p{(\columnwidth - 8\tabcolsep) * \real{0.21}}
  >{\raggedleft\arraybackslash}p{(\columnwidth - 8\tabcolsep) * \real{0.21}}
  >{\raggedleft\arraybackslash}p{(\columnwidth - 8\tabcolsep) * \real{0.21}}@{}}
\toprule
Nano configuration & Magic RMSE & LICS mIoU & MARIDA mIoU & MADOS
mIoU \\
\midrule
\endhead
Official Nano & 0.633355 & 0.955118 & \textbf{0.647160} &
\textbf{0.565515} \\
Coastal adaptation, no MoE & 0.647476 & \textbf{0.969692} & 0.507578 &
0.518878 \\
Unfrozen, physical coordinates, Top-1 & 0.633957 & 0.961896 & 0.486897 &
0.515103 \\
Unfrozen, physical coordinates, Top-2 & 0.632331 & 0.962021 & 0.491778 &
0.522514 \\
Frozen, physical coordinates, Top-1, epoch 40 & 0.568042 & 0.935636 &
0.540643 & 0.492267 \\
Frozen, physical coordinates, Top-2, epoch 40 & \textbf{0.564655} &
0.925591 & 0.537782 & 0.498283 \\
\bottomrule
\end{longtable}

Table 4. Coastal adaptation and transfer. Frozen/unfrozen refers to the
official encoder during coastal training; all feature extractors are
frozen during downstream probing. The no-MoE encoder contains
approximately 1.700 million parameters and the Nano MoE encoders
approximately 2.274 million, excluding downstream heads.

\hypertarget{optical-and-depth-foundation-features-are-complementary-on-two-marine-tasks}{%
\subsection{4.3. Optical and depth foundation features are complementary
on two marine
tasks}\label{optical-and-depth-foundation-features-are-complementary-on-two-marine-tasks}}

The frozen OLMoEarth--PromptDA dual-foundation representation improves
MARIDA by 2.55 percentage points and MADOS by 3.54 points relative to
OLMoEarth Large alone (Table 5). Bathymetry RMSE changes from 0.577098
to 0.574892 m, while LICS decreases by 1.16 points. The dual-foundation
score remains below the best reference in the full model survey on each
task, including Nano on MARIDA and TerraMind on MADOS. The result
nevertheless supports the specific method claim: features from a depth
foundation model can complement an optical foundation model for selected
marine transfers without target-task adaptation of either encoder.

\begin{longtable}[]{@{}
  >{\raggedright\arraybackslash}p{(\columnwidth - 10\tabcolsep) * \real{0.14}}
  >{\raggedright\arraybackslash}p{(\columnwidth - 10\tabcolsep) * \real{0.14}}
  >{\raggedleft\arraybackslash}p{(\columnwidth - 10\tabcolsep) * \real{0.18}}
  >{\raggedleft\arraybackslash}p{(\columnwidth - 10\tabcolsep) * \real{0.18}}
  >{\raggedleft\arraybackslash}p{(\columnwidth - 10\tabcolsep) * \real{0.18}}
  >{\raggedleft\arraybackslash}p{(\columnwidth - 10\tabcolsep) * \real{0.18}}@{}}
\toprule
Large-based configuration & Before probing & Magic RMSE & LICS mIoU &
MARIDA mIoU & MADOS mIoU \\
\midrule
\endhead
Official OLMoEarth Large & No adaptation & 0.577098 & 0.954755 &
0.612900 & 0.647969 \\
OLMoEarth Large + PromptDA depth foundation & No adaptation & 0.574892 &
0.943147 & 0.638368 & 0.683345 \\
Large + PromptDA + SpectralMoE & 20 epochs on each target task &
0.560616 & 0.954343 & 0.495718 & 0.594606 \\
\bottomrule
\end{longtable}

Table 5. OLMoEarth-only, frozen optical--depth fusion, and
target-task-adapted SpectralMoE settings. The last row has additional
target-task supervision and uses patch eight rather than patch four. It
is not a controlled test of MoE versus the dual-foundation method. The
frozen dual-foundation combination contains approximately 612.139
million parameters; the task-adapted MoE combination contains
approximately 617.877 million, excluding unused depth prediction and
downstream heads.

The task-adapted MoE configuration does not recover the MARIDA or MADOS
performance of the frozen optical--depth fusion, despite having received
additional labels before probing. Differences in patch size, interface,
head, and optimization preclude a causal claim that MoE is inferior.
They do show why a dual-foundation method should first be compared with
each frozen encoder and why target-task adaptation should not be
conflated with complementary pretrained features.

\hypertarget{discussion-preserving-representations-while-learning-coastal-corrections}{%
\section{5. Discussion: preserving representations while learning
coastal
corrections}\label{discussion-preserving-representations-while-learning-coastal-corrections}}

Together, the experiments support a focused interpretation of the
dual-foundation method. OLMoEarth supplies optical and multispectral
evidence, whereas PromptDA supplies a depth-oriented visual
representation. Their frozen combination improves two marine-surface
transfers, but the unchanged encoders do not guarantee gains on every
task. The method therefore offers a low-supervision baseline for testing
whether geometric priors learned outside remote sensing can complement
optical Earth-observation features.

A concrete next experiment is to learn a gated correction on top of the
dual-foundation representation while preserving both original feature
paths:

\[
F_{\mathrm{out}}=F_0+\alpha(x)\,\Delta(F_0,F_{\mathrm{aux}}).
\]

Here, both pretrained encoders can remain frozen, the residual output
can start at zero, and a feature-preservation penalty can discourage
unnecessary departures from the two pretrained representations. This is
a proposed experiment, not a method validated in this paper, and
residual gating and distillation are established design ideas rather
than claimed inventions. Its testable question is whether coastal
supervision can improve the dual-foundation representation without
removing the MARIDA and MADOS gains observed under frozen probing. The
minimum controls are OLMoEarth alone, PromptDA alone, frozen
concatenation, a single-encoder adapter and a two-encoder adapter with
matched training data and search budgets.

Several boundaries affect interpretation. First, the reported
configurations are single runs without confidence intervals; small
differences, especially the 0.38\% bathymetry improvement from
concatenation, should not be treated as established gains. Second,
changes in data normalization, sensor mapping, and feature dimensions
complicate cross-model ranking. Third, the coastal subset is
geographically imbalanced and spatial overlap with evaluation data has
not been audited. Fourth, PromptDA contributes features learned in a
different domain; neither the use of its encoder nor improved scores
proves that metric depth is the relevant information. Finally, four
optical tasks do not cover sea-surface temperature, chlorophyll,
sediment concentration, SAR performance, or full-year temporal
reasoning.

These boundaries suggest focused follow-up tests. A dimension-matched
fusion control would test whether extra probe capacity explains part of
the observed benefit. A learned projection or low-rank fusion would test
whether channel concatenation is sufficient. Matched-seed and
fixed-patch comparisons would separate preprocessing from encoder
complementarity. Paired Sentinel-1 inputs could then assess whether the
optical--depth combination remains useful when a second remote-sensing
modality is added. None of these outcomes is assumed by the present
study.

\hypertarget{conclusion}{%
\section{6. Conclusion}\label{conclusion}}

In the evaluated coastal setting, combining OLMoEarth with a frozen
depth foundation model provides complementary features for MARIDA and
MADOS under a common linear-probing protocol. The same dual-foundation
representation does not improve every task, and its larger probe and
interface differences limit causal attribution. The result establishes a
concrete baseline for optical--depth fusion in coastal remote sensing
and motivates future learned fusion that preserves both pretrained
feature paths.

\hypertarget{data-code-and-manuscript-status}{%
\section{Data, code, and manuscript
status}\label{data-code-and-manuscript-status}}

The project repository is
\url{https://github.com/NNNoreply/Coastal_Foundation_Models}. Its
initial release contains a project overview; code, model checkpoints,
and data are not yet included. This working draft is based on existing
local experiment records. Author names, affiliations, contact details,
dataset permissions, and complete bibliographic metadata remain to be
supplied before submission. Existing model and dataset creators retain
attribution for their resources. AI assistance was used to organize and
draft this manuscript from local code and result records; author
verification of the submitted text and interpretations remains pending.

\hypertarget{source-references}{%
\section{Source references}\label{source-references}}

The following are source pointers verified in the local repository
documentation, not a completed bibliography. Canonical citations for
every benchmark, model, geospatial product, optimizer, and proposed
residual/distillation context remain to be added and checked before
public submission.

\begin{enumerate}
\def\labelenumi{\arabic{enumi}.}
\tightlist
\item
  OLMoEarth project. Official model and evaluation implementation.
  \url{https://github.com/allenai/olmoearth_pretrain}. The repository
  links the original report at \url{https://arxiv.org/abs/2511.13655}.
  Nano and Large in this study use different released model versions.
\item
  Xi Chen, Maojun Zhang, Yu Liu, and Shen Yan. \emph{Local Precise
  Refinement: A Dual-Gated Mixture-of-Experts for Enhancing Foundation
  Model Generalization against Spectral Shifts}. arXiv:2603.13352, 2026.
  \url{https://arxiv.org/abs/2603.13352}. Implementation:
  \url{https://github.com/daxichen/SpectralMoE-main}. Author/title
  metadata follow the released repository citation.
\item
  Haotong Lin et al.~\emph{Prompting Depth Anything for 4K Resolution
  Accurate Metric Depth Estimation}. arXiv, 2024.
  \url{https://promptda.github.io/}. The experiment uses the released
  PromptDA ViT-L visual encoder checkpoint.
\item
  \emph{CROMA: Remote Sensing Representations with Contrastive
  Radar-Optical Masked Autoencoders}. Paper linked in the local official
  implementation: \url{https://arxiv.org/abs/2311.00566}.
\item
  SatMAE project and paper, linked in the local official implementation:
  \url{https://sustainlab-group.github.io/SatMAE/} and
  \url{https://arxiv.org/abs/2207.08051}.
\item
  \emph{Scale-MAE: A Scale-Aware Masked Autoencoder for Multiscale
  Geospatial Representation Learning}. Paper linked in the local
  implementation: \url{https://arxiv.org/abs/2212.14532}.
\end{enumerate}

\end{document}
